\cotefile{157-1}
\scanpage{18}
\mspage{15}
% archive sheet number: 17
\margin{\deleted{\unsure{18}}}
Par exemple, considérons le dérivateur
\[
\mathrm{HOTAB}
\]
dont la cat. \unsure{fondamentale} est $\mathrm{Hotab}$, la \inserted{sous-cat. pleine de la} cat. dérivée
\deleted{des complexes de \uline{chaînes}} \inserted{$\mathrm D(\mathrm{Ab})$} \deleted{de} $\mathrm{Ab}$, \deleted{\unsure{soit}} formée des complexes \deleted{\unsure{qui sont}} à coh. nulle en degrés $>0$.
On peut décrire $\mathrm{HOTAB}$ par \deleted{\unsure{la|les}} catégorie de modèles $\underline{\mathrm{Hom}}(\unsure{\underline{\Delta}^{\circ}|\Delta^{\circ}},\mathrm{Ab})$ des groupes abéliens $1/2$ simpliciaux, avec comme équivalences \deleted{faibles} les quasi-isom. au sens habituel. Le ``\uline{foncteur d'abélianisation}'' \deleted{\unsure{de passage au}} \inserted{classique}
\[
\mathrm{Hot}\longrightarrow\mathrm{Hotab}
\]
se prolonge en un morphisme de dérivateurs
\[
\mathrm{HOT}\longrightarrow\mathrm{HOTAB}
\]
La question de sa compatibilité avec $f_*$ ne semble pas tautologique. Elle \deleted{\unsure{revient}} impliquerait p.ex. que si
\diagram{\begin{tikzcd}[ampersand replacement=\&]
X \arrow[d,"f"'] \& X' \arrow[l,"g"'] \arrow[d,"f'"]\\
Y \& Y' \arrow[l,"{\unsure{j|\delta}}"]
\end{tikzcd}}
est un diagramme cartésien dans $\mathcal M=\underline{\mathrm{Hom}}(\Delta^\circ,\mathrm{Ens})$,
avec $f$ une fibration (\unsure{en restant} dans le cadre des esp. top., \unsure{une|en} \inserted{p.ex.} fibration loc. triviale) alors le \unsure{carré} dans $\mathrm{Hotab}$
% slanted
\margin{$\unsure{\longleftrightarrow|\longleftarrow}$ C'est faux déjà si \deleted{\unsure{le}} $Y$ est l'espace ponctuel, \unsure{puisque|très} \unsure{visiblement} le $\mathbb R\Gamma_!$ ne transforme pas produit en \unsure{une|en|ou} \unsure{somme} ! Ce foncteur \unsure{n'est guère|n'est pas un} ``exact à \unsure{gauche|g.}'' (comme morphisme de dérivateurs)}
\diagram{\begin{tikzcd}[ampersand replacement=\&]
\mathbb R\Gamma_!(X)\arrow[d] \& \mathbb R\Gamma_!(X')\arrow[l]\arrow[d]\\
\mathbb R\Gamma_!(Y) \& \mathbb R\Gamma_!(Y')\arrow[l]
\end{tikzcd}}
est \inserted{hot-}cartésien. \deleted{\illegible{2}} Et de $\unsure{=|\Leftarrow}$, \unsure{pour|par} une hot-\deleted{\unsure{so}}\unsure{somme amalgamée} (carré \inserted{hot-}cocartésien dans $\mathrm{Hot}$).
